\chapter{Naming the Field}

What we're describing lacks a settled name. The practice of effective human-AI cognitive collaboration is too new to have established terminology. Yet naming matters: it enables curriculum development, professional identity, and systematic improvement.

\section{The Emergence of New Terms}

The field is rapidly developing vocabulary. Since 2024, several terms have gained traction:

\subsection{Prompt-Driven Development (PDD)}

Prompt-Driven Development has emerged as a term for software development centered on AI dialogue rather than traditional coding:

\begin{quotebox}
``Prompt-Driven Development represents a fundamental shift in how software is conceived, built, and maintained. Rather than manually writing every function, developers can now focus on articulating goals, solving problems, and curating the outputs of AI collaborators.''
\end{quotebox}

PDD focuses on the \emph{what}---using prompts to direct AI to produce artifacts. DCF complements this by focusing on \emph{how to think} during the process.

\subsection{Vibe Coding}

In 2025, ``vibe coding'' emerged as a colloquial term for conversational, prompt-first development workflows:

\begin{itemize}
    \item Developers describe what they want in natural language
    \item AI generates working code in response
    \item The emphasis is on rapid iteration over precision
\end{itemize}

The term captures the intuitive, exploratory nature of AI-assisted development but lacks rigor. Current applications focus on prototyping and MVPs rather than production systems.

\subsection{AI-Driven Development Lifecycle (AI-DLC)}

AWS introduced this term in 2025 for an AI-native methodology placing AI at the center of the software development process:

\begin{itemize}
    \item AI assists at every stage from inception to operation
    \item Human oversight remains at every step
    \item Focus on effectively applying AI at each development phase
\end{itemize}

\subsection{Spec-Driven Development}

A response to the chaos of pure prompt-driven approaches, spec-driven development uses formal specifications as ``super prompts'':

\begin{itemize}
    \item Specifications are version-controlled and human-readable
    \item They express programmer intent crisply to agents
    \item Structure tames the ad-hoc nature of prompting
\end{itemize}

\section{Related Established Terms}

Several existing terms overlap with or inform the emerging discipline:

\begin{table}[H]
\centering
\begin{tabular}{lp{7cm}}
\toprule
\textbf{Term} & \textbf{Focus} \\
\midrule
Augmented Intelligence & Human capability extended by AI (contrast with ``artificial'' replacement) \\
Human-AI Teaming & Collaboration patterns between humans and AI systems \\
Prompt Engineering & Crafting effective prompts (narrower than cognitive methodology) \\
Cognitive Augmentation & Enhancement of human thinking through technology \\
\bottomrule
\end{tabular}
\caption{Related Established Terms}
\end{table}

\section{Candidate Names for the Discipline}

Drawing on the emerging terminology and established fields, several candidates emerge for naming what DCF addresses:

\begin{itemize}
    \item \textbf{Co-cognitive engineering}: Building systems of thought with AI---emphasizes the collaborative construction aspect
    \item \textbf{Augmented intelligence practice}: Human capability extended by AI---emphasizes augmentation over replacement
    \item \textbf{Knowledge architecture}: Designing how individuals and organizations think---emphasizes the structural aspect
    \item \textbf{Cognitive systems design}: Engineering human-AI thought systems---emphasizes the systemic nature
    \item \textbf{AI collaboration methodology}: How to work effectively with AI---emphasizes the methodological rigor
\end{itemize}

\section{Why Naming Matters}

The lack of a settled name isn't merely semantic. It has practical consequences:

\subsection{Curriculum Development}

Without a name, you can't create courses. ``Introduction to [what]?'' Universities and training programs need terms to organize offerings.

\subsection{Professional Identity}

Professionals need language to describe their expertise. ``I'm skilled at effective AI collaboration'' lacks the precision of a named discipline.

\subsection{Research Coordination}

Academic research requires shared vocabulary. Without it, researchers talk past each other, duplicate work, and miss connections.

\subsection{Organizational Recognition}

Companies struggle to hire for roles they can't name. ``Prompt engineer'' captures only part of what effective practitioners do.

\section{What DCF Contributes}

Regardless of what the broader field is called, DCF contributes a specific focus:

\begin{quotebox}
\textbf{Most frameworks focus on what the AI should do. DCF focuses on how the human should think.}
\end{quotebox}

\begin{table}[H]
\centering
\begin{tabular}{ll}
\toprule
\textbf{Other Approaches} & \textbf{DCF} \\
\midrule
How to prompt effectively & How to think effectively with AI \\
What outputs to generate & What questions to ask \\
AI capabilities & Human cognitive stance \\
Task completion & Cognitive development \\
\bottomrule
\end{tabular}
\caption{DCF's Distinctive Focus}
\end{table}

The field needs both: methodologies for directing AI \emph{and} frameworks for human thinking during collaboration. DCF addresses the latter.

\section{The Recognition That Matters}

Beyond any specific name, the crucial recognition is this:

\begin{quotebox}
\textbf{This is a discipline.}
\end{quotebox}

Whether called co-cognitive engineering, AI collaboration methodology, or something not yet coined, the practice of thinking effectively with AI:

\begin{itemize}
    \item Can be learned
    \item Can be practiced systematically
    \item Can be improved over time
    \item Can be taught to others
    \item Can be mastered
\end{itemize}

The name will settle. What matters now is treating this as a real skill domain deserving deliberate development.

